\documentclass[10pt,a4paper]{amsart}
\usepackage[margin=19mm]{geometry}
\usepackage{amsmath,amssymb,mathtools}
\usepackage[scr=rsfs,cal=euler]{mathalfa}
\usepackage{xfrac}
\usepackage[unicode,hidelinks,bookmarks=false]{hyperref}
\newcommand{\ie}{i.e.\ }
\newcommand{\cf}{cf.\ }
\newcommand{\resp}{respectively\ }
\newcommand{\Subsection}{\subsection}
\providecommand{\nobreakdash}{\nobreak\hbox{-}}
\setlength{\emergencystretch}{2em}
\multlinegap0pt
\usepackage{xcolor}
\usepackage{enumerate}
\usepackage{mathtools}
\usepackage{float}
\usepackage{tikz}
\newtheorem{lemma}{Lemma}
\newcommand{\E}{\mathrm{E}}
\newcommand{\combi}[2]{\begin{pmatrix}#1 \\ #2 \end{pmatrix}}
\title{Limiting partition function for the Mallows model: a conjecture and partial evidence}
\author{Soumik Pal}
\date{Source: arXiv:2406.18855v2 (CC BY 4.0)}
\begin{document}
\maketitle

\noindent\emph{Source and licence.} Limiting partition function for the Mallows model: a conjecture and partial evidence. Version arXiv:2406.18855v2 (CC BY 4.0), distributed by arXiv under the Creative Commons Attribution 4.0 International licence, http://creativecommons.org/licenses/by/4.0/, as recorded in the arXiv licence metadata for this exact version and shown on the abstract page https://arxiv.org/abs/2406.18855v2. Full licence notice: \texttt{licenses/arxiv-2406.18855-CC-BY-4.0.txt}.\par\smallskip

\noindent\textit{Editorial scope.} Counted unit: the article's lemma lemma (source0.tex lines 497--502). The article's complete original proof of it is delivered with it (source0.tex lines 504--598, one \texttt{proof} environment of 95 lines). No mathematics is written, repaired or re-derived: the statement and the proof are the article's own lines, in the article's own notation, and no retained line is substituted or re-flowed.  Retained uncounted as the source-local context the proof needs: lines 252--254.  Nothing was omitted from any retained span, so the package carries no omission record.  No retained line contains a citation, so the wrapper carries no bibliography and no bibliography entry.  No retained line contains a figure environment, an image-inclusion command or drawing code, so no image is delivered and the package declares no asset dependency.  Editorial formatting only, in addition: an AMS \texttt{amsart} preamble that restates the article's own declarations and macros used by the retained lines, namely the environments \texttt{lemma} on separate counters, together with the standard packages those lines need.  The retained environments print in this document as Lemma 1 in order of appearance, whereas the article prints the same items with the numbering of its own full document.  The complete native source of the article as distributed in the arXiv e-print for this exact version is bundled unchanged as \texttt{sources/161438/source0.tex}.
\begin{equation}\label{eqn:marginal}
    \int_0^1 \tilde{\rho}(x,y)dy=0=\int_0^1 \tilde{\rho}(z,w)dz, 
\end{equation}

\begin{lemma}
There is a $C_0>0$ and $\sigma \in (0,1)$ such that 
\[
\E\left[ \prod_{t=1}^r \tilde{\rho}(U_{i}, V_i) \mid E_{n,r} \cap F_{n,r} \right] \le C_0 \frac{\sigma^r}{n^r}. 
\]
\end{lemma}

\begin{proof}
  \textcolor{red}{assume all integrals are with respect to continuous uniform.}  For $r=1$, $\E(\tilde{\rho}(U_1, V_1) \mid U_1)=0$, due to \eqref{eqn:marginal}. Hence so is the full expectation. 

  For $r=2$, on the event $E_{n,r} \cap F_{n,r}$,
  \[
  \E\left[ \tilde{\rho}(U_2, V_2) \mid U_1, V_1, U_2 \right]= \frac{1}{n} \sum_{j\neq V_1} \tilde{\rho}(U_2, j)= - \frac{1}{n}\tilde{\rho}(U_2, V_1),
  \]
  by the conditional mean-zero property \eqref{eqn:marginal}. Similarly, on the event $E_{n,r} \cap F_{n,r}$,
  \[
  -\E\left[ \tilde{\rho}(U_2, V_1) \mid U_1, V_1 \right]=  \frac{1}{n}\tilde{\rho}(U_1, V_1).
  \]
  Thus
  \[
  \E\left[ \tilde{\rho}(U_1, V_1)\tilde{\rho}(U_2, V_2)
  \right]= \frac{1}{n^2} \E\left[ \tilde{\rho}^2(U_1, V_1)\right]\le \frac{\sigma^2}{n^2}.
  \]

The case for a general $r$ can now be done by iterating the above argument. Given $\prod_{i=1}^r \tilde{\rho}(U_i, V_i)$, we will progressively compute the conditional expectations given the backward filtration for $k=1,2,\ldots, r-1$, 
\[
\begin{split}
\mathcal{G}_{2k-1}&=\sigma\left( U_1, \ldots, U_{r-(k-1)}, V_1, \ldots, V_{r-k}\right)\\
\mathcal{G}_{2k}&=\sigma\left( U_1, \ldots, U_{r-k}, V_1, \ldots, V_{r-k}\right).
\end{split}
\]
The conditional expectation will be computed given the event $E_{n,r} \cap F_{n,r}$.  

Thus, from our computation for $r=2$,  
\[
\begin{split}
\E&\left[ \prod_{i=1}^r \tilde{\rho}(U_i, V_i) \mid \mathcal{G}_1 \right]=-\prod_{i=1}^{r-1} \tilde{\rho}(U_i, V_i) \left( \frac{1}{n} \sum_{i=1}^{r-1} \tilde{\rho}(U_r, V_i) \right)\\
\E&\left[ \prod_{i=1}^r \tilde{\rho}(U_i, V_i) \mid \mathcal{G}_2 \right]=\prod_{i=1}^{r-1} \tilde{\rho}(U_i, V_i) \frac{1}{n^2} \sum_{i=1}^{r-1} \sum_{j=1}^{r-1} \tilde{\rho}(U_i, V_j)\\
&= \frac{1}{n^2} \sum_{i=1}^{r-1} \tilde{\rho}^2(U_i, V_i) \prod_{k\neq i,\; k\in [r-1]} \tilde{\rho}(U_i, V_i)  \\
&+ \prod_{i=1}^{r-1} \tilde{\rho}(U_i, V_i) \frac{1}{n^2} \sum_{i\neq j} \tilde{\rho}(U_i, V_j).
\end{split}
\]
Hence, by exchangeability, 
\[
\begin{split}
\E&\left[ \prod_{i=1}^r \tilde{\rho}(U_i, V_i) \mid E_{n,r} \cap F_{n,r} \right]= \frac{(r-1)}{n^2} \E\left[ \tilde{\rho}^2(U_1, V_1) \prod_{k=2}^{r-1} \tilde{\rho}(U_i, V_i) \mid E_{n,r} \cap F_{n,r} \right]\\
&+ \frac{1}{n^2} \combi{n}{2} \E\left[ \tilde{\rho}(U_1, V_2)\prod_{k=1}^{r-1} \tilde{\rho}(U_i, V_i) \mid E_{n,r} \cap F_{n,r} \right].
\end{split}
\]

This evolving terms may be represented by a sequence of (multi)graphs on vertices $[r]:=\{1,2,\ldots, r\}$ and $[r']:=\{1', 2', \ldots, r'\}$. The graph always has $r$ edges although multiple edges are allowed. Initially the set of edges are given by $\{(i,i'),\; i \in [r]\}$. Successively, every vertex of degree one is removed, the edge incident on that vertex is deleted, and an edge is added between the remaining vertices. This is continued until no vertex of degree one remains. Every time a vertex is removed a factor of $-n^{-1}$ gets multiplied. Clearly, we remove the minimum number of vertices when at the end each remaining vertex has degree two. Since the number of edges (counting multiplicities) always remains the same, at least $2\lceil r/2 \rceil$ many vertices will be removed and the final term will contribute a factor of $n^{-2\lceil r/2 \rceil}$ or higher powers of $n^{-1}$. 

The terms that contain exactly the factor $n^{-2\lceil r/2 \rceil}$ are those where every remaining $2\lfloor r/2 \rfloor$ many vertices has degree exactly two. The number of such graphs are exactly the number of permutations of $[\lfloor r/2 \rfloor]$ which is $\lfloor r/2 \rfloor !$. If the terminal graph has $k$ many edges with incident vertices whose degree is more than two, it comes with a factor of $n^{-2\lceil r/2 \rceil + k}$, for $k=1, \ldots, r-1$. 

\newpage

\textbf{A new approach.}

Consider the term 
\[
\frac{(n-r)!}{r! n!} n^{2r} \E\left[ \prod_{t=1}^r \tilde{\rho}(U_i, V_i) 1\{ E^{n,r} \cap F^{n,r}\} \right].
\]
We will apply the Inclusion-Exclusion principle for each of these terms. Similar to before, let's define events 
\[
E_{ij}= \{U_i=U_j\}, \quad F_{ij}=\{V_i = V_j\}, \quad 1\le i < j \le r.
\]
Above and below we are going to ignore the difference between integration with respect to the  continuous and discrete uniform measures. 
\[
\begin{split}
    \E&\left[ \prod_{t=1}^r \tilde{\rho}(U_i, V_i) 1\{ E^{n,r} \cap F^{n,r}\} \right]= \E\left[ \prod_{t=1}^r \tilde{\rho}(U_i, V_i) \right]\\
    &- \E\left[ \prod_{t=1}^r \tilde{\rho}(U_i, V_i); \cup_{1\le i < j \le r} E_{ij} \cup F_{ij} \right]\\
    &= - \E\left[ \prod_{t=1}^r \tilde{\rho}(U_i, V_i); \cup_{1\le i < j \le r} E_{ij} \cup F_{ij} \right]=\sum_{k,l=0,\; k+l>0}^{r(r-1)/2} (-1)^{k+l-1} \\
    \times &\sum \E\left[ \prod_{t=1}^r \tilde{\rho}(U_i, V_i); \text{intersection of $k$ many $E_{ij}$ and $l$ many $F_{ij}$}  \right].
\end{split}
\]
The inner sum is over all choices of $k$ many $E_{ij}$ and $l$ many $F_{ij}$. Fix $(i_1, j_1), \ldots (i_k, j_k)$ and $(i'_1, j'_1), \ldots, (i'_l, j'_l))$. For every such choice the event $\cap_{m=1}^k E_{i_mj_m}$ can be represented by a graph with vertices $[r]$ and $i_m \sim j_m$ for $m\in [k]$. This graph partitions the set $[r]$ in blocks given by its connected subsets. The event then means that all vertices in the same connected component must have the same $U$ values. Thus $P(\cap_{m=1}^k E_{i_mj_m})= n^{-r + \#\text{blocks}}$. A similar representation holds for $\cap_{m=1}^l F_{i'_mj'_m}$. Thus 
\[
P\left(\cap_{m=1}^k E_{i_mj_m} \cap \cap_{m=1}^l F_{i'_mj'_m} \right)= n^{-2r + \sum \#\text{blocks}},
\]
where the sum is over the number of blocks in the two partitions. 





\[
\begin{split}
\E&\left[ \prod_{t=1}^r \tilde{\rho}(U_i, V_i); \text{intersection of $k$ many $E_{ij}$ and $l$ many $F_{ij}$}  \right]\\
=&\E\left[ \prod_{t=1}^r \tilde{\rho}(U_i, V_i);  \cap_{m=1}^k E_{i_mj_m} \cap_{m=1}^l F_{i'_mj'_m} \right] \\
=& \E\left[ \prod_{t=1}^r \tilde{\rho}(U_i, V_i) \mid   \cap_{m=1}^k E_{i_mj_m} \cap_{m=1}^l F_{i_mj_m} \right] P\left( \cap_{m=1}^k E_{i_mj_m} \cap_{m=1}^l F_{i'_mj'_m} \right)\\
=& \E\left[ \prod_{t=1}^r \tilde{\rho}(U_i, V_i) \mid   \cap_{m=1}^k E_{i_mj_m} \cap_{m=1}^l F_{i'_mj'_m} \right] n^{-2r + \sum \#\text{blocks}}.
\end{split}
\]

Consider now the conditional expectations. Every block in the partition takes the same value as a single uniform random variable. Disjoint blocks take independent values. Thus, if there is any block of size one, the entire term will vanish. Hence, the number of blocks is at least $r/2$ to have a nonzero contribution. Then,
\[
\frac{(n-r)!}{r! n!} n^{2r-p-q}= \frac{(n-r)!}{r! n!} n^{\#\text{cycles}+ \#\text{cycles}'}.
\]
Hence, in the worst possible case of all cycles of size two, we get a factor of $n^r$. In all other cases, we get a smaller order factor. 

  
\end{proof}

\end{document}
